\scanpage{82}
Inversement, soit $P_1$ \inserted{\unsure{une}} catégorie \inserted{$\otimes$}\unsure{-}additive \inserted{avec},
munie d'un \inserted{$\otimes$ \deleted{\unsure{ass}} \unsure{avec unité|avec objet unité|ACU strict}} foncteur de $(P_1,+)$
dans $(P_{1+}^{\mathrm{grad.}},\otimes)$, \deleted{\unsure{deux} \unsure{notif}}
\[
\Gamma^*:(P_1,\oplus)\longrightarrow(P_1^{\mathrm{grad}},\otimes)
\]
donc
\[
\Gamma^*(0)=\mathbb{1}\qquad\text{i.e.}\qquad
\Gamma^i(0)=\begin{cases}\mathbb{1}&\text{si }i=0\\0&\text{si }i>0\end{cases}
\]
\[
\Gamma^*(\deleted{M}X\oplus\deleted{N}Y)\simeq
\overbrace{\Gamma^*(X)\otimes\Gamma^*(\unsure{M|Y})}^{\text{avec degr \unsure{tot.|total}}}
\]
(isom \deleted{\unsure{de gr.|gr.}} donné, bifonctoriel en $X,Y$, compatible
avec contraintes d'unité, d'ass., de comm.)
avec
\[
\Gamma^i(M)=\begin{cases}0&\text{si }i<0\\\mathbb{1}&\text{si }i=0\\M&\text{si }i=1\end{cases}
\]
(isom de foncteurs en $M$), Supposons de
plus donné des hom de foncteurs
\[
\Gamma^{pq}\xrightarrow{\alpha_{pq}}\Gamma^p\deleted{\unsure{(}}\circ\Gamma^q
\]
(i.e.
\[
\Gamma^{pq}(X)\xrightarrow{\alpha_{pq}^X}\Gamma^p(\Gamma^q(X))
\]
)
satisfaisant la condition d'associativité
ci dessus, plus les conditions
évidentes pour \deleted{$p,q$} l'un des nbs $p,q$
égaux à $0$ ou $1$ :
\diagram{\begin{tikzcd}[ampersand replacement=\&,column sep=large]
{(**)} \& \Gamma^q\arrow[r,"{\alpha_{1,q}=\mathrm{id}}"]\arrow[dr,bend right=24,"{\alpha_{q,1}=\mathrm{id}}"'] \& {\Gamma^1\circ\Gamma^q=\mathrm{id}\circ\Gamma^q=\Gamma^q}\\
\& \& {\Gamma^q\circ\Gamma^1=\Gamma^1\circ\mathrm{id}=\Gamma^q}
\end{tikzcd}}
\diagram{\begin{tikzcd}[ampersand replacement=\&,column sep=large]
{(***)} \& \Gamma^0\arrow[r,"{\alpha_{0,q}=\mathrm{id}}"]\arrow[dr,bend right=24,"{\alpha_{q,0}\ \text{un }\text{\uuline{iso}}}"'] \& {\Gamma^0\circ\Gamma^q=\Gamma^0}\\
\& \& {\Gamma^q\circ\Gamma^0.}
\end{tikzcd}}
% archive sheet number: 81
